\begin{claim}{D.44}{Proved}{A first-symmetry choice for the involution coboundary}
Let $\alpha=(-1,f)\in GT_0$ be an involution in its
normalized actual action on $F=\widehat F(x,y)$:
\[
 \alpha(x)=x^{-1},\qquad
 \alpha(y)=f^{-1}y^{-1}f,\qquad f\in F'.
\]
Write $\theta(x)=y$, $\theta(y)=x$.
There exists $h\in F'$ satisfying both exact equations
\[
               f=h^{-1}\alpha(h),\qquad h\theta(h)=1.
                                                        \tag{D.44a}
\]
Thus the coboundary in D.35 can be chosen with the
first Harbater--Schneps symmetry.  For every unit
$\mu\in\widehat\Z^\times$, the associated continuous
endomorphism
$\rho_{\mu,h}(x)=x^\mu$,
$\rho_{\mu,h}(y)=hy^\mu h^{-1}$ satisfies
$\alpha\rho_{\mu,h}=\rho_{\mu,h}\delta$, where
$\delta$ inverts $x,y$.
Its full peripheral-conjugator condition is precisely
D.35's remaining equation
$x^\mu hy^\mu h^{-1}\sim_F(xy)^\mu$.

At the characteristic finite quotient $A_5^{19}$
of D.38, this first-symmetry condition is soluble
for every one of the $5\,062\,500$ tuples that
pass the triangle test.  This finite equation
therefore gives no further exclusion among them.
\end{claim}
\begin{proof}
The normalized involution and the first
Harbater--Schneps equation give
\[
                  f\alpha(f)=1,\qquad\theta(f)f=1.
                                                        \tag{D.44b}
\]
Checking the two free generators gives the
identity of actual automorphisms
\[
                 \theta\alpha\theta
                        =\operatorname{Inn}(f)\alpha.
                                                        \tag{D.44c}
\]
Indeed both sides send $x$ to $fx^{-1}f^{-1}$
and $y$ to $y^{-1}$.

Form $E=F\rtimes_\alpha\langle s\mid s^2=1\rangle$.
Define an automorphism $\Theta$ of this exact
extension by
\[
                      \Theta|_F=\theta,\qquad
                      \Theta(s)=fs.                    \tag{D.44d}
\]
The first equation of (D.44b) preserves $s^2=1$;
(D.44c) preserves the conjugation relation;
and the second equation of (D.44b) gives
$\Theta^2(s)=\theta(f)fs=s$.
Thus $\Theta^2=1$ on all of $E$.
Adjoin its named involution $\tau$ and put
\[
 G=E\rtimes_\Theta\langle\tau\rangle,\qquad
                         G/F=C_2\times C_2.
                                                        \tag{D.44e}
\]
The two quotient coordinates are represented
by $s,\tau$.

The preimage of the $\tau$ line is the normal
subgroup
$G_\theta=F\rtimes_\theta\langle\tau\rangle$.
It has the explicit free-product identification
\[
 G_\theta\xrightarrow{\sim}
                  C_2(\tau)\amalg\widehat\Z(v),
       \qquad x\mapsto v,\quad y\mapsto\tau v\tau.
                                                        \tag{D.44f}
\]
The inverse sends $v$ to $x$ and $\tau$ to
the named involution; both composites fix
the indicated generators.
The free-factor intersection theorem
\cite[Theorem 5 of the inspected author version]{A-HerfortZalesskii2008}
therefore gives $C_F(\tau)=1$.

There is exactly one $F$-conjugacy class of
involutions in $G_\theta$.
For a direct verification, its mod-two
$H^1(F,\F_2)$ module is the regular
two-dimensional $C_2$ module, since $\theta$
exchanges its two basis characters.
Its positive-degree $C_2$ cohomology vanishes.
The split Hochschild--Serre calculation
therefore gives
$\dim H^n(G_\theta,\F_2)=1$ for $n\geq3$.
Brown restriction
\cite[Theorem 2 and \S1.1]{A-KochloukovaZalesskii2007}
gives one involution class, by exactly the
dimension argument of (D.37f).
Conjugacy by $G_\theta$ equals conjugacy by
$F$, since $G_\theta=F\langle u\rangle$
for every involution $u$ and its second
factor centralizes that element.

Consider the closed normalizer
$V=N_G(\langle\tau\rangle)$.
For every $g\in G$, the conjugate $g\tau g^{-1}$
belongs to $G_\theta$, and its image in
$G/F$ is still the $\tau$ coordinate.
The preceding uniqueness gives an $F$-conjugator
taking it back to $\tau$.
It follows that $V\to G/F$ is onto.
Its kernel is $C_F(\tau)=1$, so
\[
                   V\xrightarrow{\sim}C_2\times C_2.
                                                        \tag{D.44g}
\]
In particular $V$ is a finite subgroup,
contains $\tau$, and has a unique involution
$u$ mapping to the $s$ coordinate.
That involution belongs to $E$ and commutes
with $\tau$.

The three $F$-conjugacy classes of involutions
in $E$ were determined in D.34.  They are
represented by
\[
                            s,\quad xs,\quad yfs.
                                                        \tag{D.44h}
\]
Their distinct mod-two abelianization vectors
also identify the classes.
Formula (D.44d) sends $s$ to $fs$, which
belongs to the class of $s$, and gives the
exact transposition
\[
                      \Theta(xs)=yfs,\qquad
                      \Theta(yfs)=xs.
                                                        \tag{D.44i}
\]
Thus the class of $s$ is the unique fixed
class under $\Theta$.
Since $u$ commutes with $\tau$, its class
is fixed, so $u=ksk^{-1}$ for some $k\in F$.
Conjugate the finite subgroup (D.44g)
by $k^{-1}$.
It now contains $s$ and
\[
                k^{-1}\tau k=h\tau,\qquad
                              h=k^{-1}\theta(k).
                                                        \tag{D.44j}
\]
The square of this second involution gives
$h\theta(h)=1$.
Its commutation with $s$ gives
$hf=\alpha(h)$, which is precisely the
first equation of (D.44a).
On abelianization, this equation reads
$-2h^{\rm ab}=f^{\rm ab}=0$.
Multiplication by two is injective on
$\widehat\Z^2$, hence $h\in F'$.
Finally substitution in $\alpha(hy^\mu h^{-1})$
using $\alpha(h)=hf$ proves
$\alpha\rho_{\mu,h}(y)=\rho_{\mu,h}(y^{-1})$; the
$x$ identity is immediate.

We check the finite scope explicitly.
Write $Q=A_5^{19}$, and use D.38's component
description of $\theta$ and its involution tuple
$t$.  The triangle test makes every component
$t_i$ a nonidentity involution of $A_5$.
The equations to solve in this quotient are
\[
 h^{-1}t h=\theta(t),\qquad h\theta(h)=1,
       \qquad f=\theta(t)t^{-1}.
                                                        \tag{D.44k}
\]
They separate over the one-element and
two-element orbits of $\theta$ on the factors.

For a two-element orbit, write
$\theta(v)_i=g v_jg^{-1}$ and
$\theta(v)_j=g^{-1}v_i g$, with $g\in S_5$.
All nonidentity involutions of $A_5$ are
conjugate within $A_5$.
Choose $h_i\in A_5$ satisfying
$h_i^{-1}t_i h_i=g t_jg^{-1}$, and set
$h_j=g^{-1}h_i^{-1}g$.
These substitutions give both equations
of (D.44k) on that orbit.

For a fixed factor, $\theta$ is conjugation
by an involution $g\in S_5$, allowing $g=1$.
Indeed its square is the identity and
$C_{S_5}(A_5)=1$.
Choose an element $v$ with $v^2=1$, commuting
with $t_i$ and having the same parity as $g$.
When $g$ is even choose $v=1$.
When $g$ is odd, take either transposition
in the disjoint-transposition expression
for $t_i$.
Then $h_i=vg^{-1}\in A_5$ satisfies both
equations in (D.44k).
Thus every triangle survivor has a finite
first-symmetry solution, including the
$2\,531\,250$ survivors failing D.38's
finite peripheral-conjugacy test.
No lifting of these finite solutions is
asserted.

The program
\path{code/part_a/pure_involution_first_symmetry.py}
executes the local construction for all
$390$ fixed-factor cases and all
$27\,000$ two-factor cases.
It checks both equations in (D.44k),
which cover every factor orbit in D.38.
\end{proof}
